% Generated TeX body; no PDF or document-style claim.
% Maintained source: problems/persistence-failure-certificates/manuscripts/correlated-fault-boundary-and-persistence-sparsifiers.md
% Recovered-source SHA-256: ebc5a8a2ab88362a202a880627c4a80de70adcac0010ee6d2d2bca78cf09cbc8
% Rebuild: python3 tools/build_manuscripts.py
% Common packages: amsmath, amssymb, hyperref, longtable, booktabs, array.
% Shared original Subset Sum preamble: manuscripts/paper-style.tex.

\hypertarget{af-3f6acd6e5b-exact-persistent-homology-certificates-under-correlated-failures}{%
\section{Exact persistent-homology certificates under correlated failures}\label{af-3f6acd6e5b-exact-persistent-homology-certificates-under-correlated-failures}}

\textbf{Abstract.} We construct exact certificates for boundary spaces and persistent homology under correlated label failures. A rank-compressed polynomial incidence lift preserves all surviving prefixes of a fixed order. The construction gives sharp retained-cell bounds and an exterior-algebra refinement for one color per vertex.

\textbf{Authors:} Yilin Liu and GPT-6 Astra Pro.

\hypertarget{af-3f6acd6e5b-main-result-and-the-quantity-being-reduced}{%
\subsection{1. Main result and the quantity being reduced}\label{af-3f6acd6e5b-main-result-and-the-quantity-being-reduced}}

Fix a coefficient field F, either the rationals or an explicitly represented finite field. Let K be a finite, explicitly given simplicial complex of dimension at most p, where p\textgreater=1. Its full (p-1)-skeleton will be retained. Let N be the number of p-simplices and let

\[
A=\partial_p(K;F),\qquad r=\mathrm{rank}_F A.
 \tag{T1}
\]

Give each vertex at most a failure labels, and give each p-simplex at most b additional labels. A query F\_0 contains at most h labels. A failed vertex and all its cofaces are removed; a p-simplex is also removed when one of its own additional labels fails. Thus a query always leaves a subcomplex. The labels may overlap arbitrarily, and a single failed label may remove arbitrarily many vertices or simplices.

There is a deterministic uniform polynomial-time algorithm retaining a set B of original p-simplices with

\[
|B|\le \min\left\{N,\ r\binom{h+a+b}{a+b}\right\}.
 \tag{T2}
\]

Let H contain B and \textbf{all original faces of dimensions at most p-1}. For any fixed input filtration (face birth precedes coface birth), the construction can ensure that the inclusions

\[
H_t-F_0\hookrightarrow K_t-F_0
 \tag{T3}
\]

induce isomorphisms in every homological degree j\textless=p-1, simultaneously for every t and every allowed F\_0. These isomorphisms commute with the filtration maps: they preserve the persistence modules, not merely individual Betti numbers. In fact the surviving p-boundary subspaces are exactly equal at every threshold.

Alternatively, given one fixed signed rational cost order on p-simplices, the same construction contains the canonical minimum-cost basis of the surviving boundary matrix after every query, including queries that decrease rank. This is a minimum-cost boundary-basis statement. Calling it a spanning acycle requires the additional homological conditions in the definition of that term.

\textbf{One order at a time.} The filtration order and an independently specified cost order cannot silently be combined in one bound (T2). A filtration-ordered construction proves the persistence assertion. A cost-ordered construction proves the cost assertion. If the two orders coincide, both hold together.

The count (T2) concerns retained \textbf{top simplices}. It is not a bound on all faces, total encoded bits, or Gaussian-elimination workspace. It does not preserve H\_p, integral torsion, or homotopy type. The selected B can depend on the chosen coefficient field.

\hypertarget{af-3f6acd6e5b-particularly-simple-case}{%
\subsubsection{Particularly simple case}\label{af-3f6acd6e5b-particularly-simple-case}}

If each vertex has one color and there are no extra top-simplex labels, then

\[
|B|\le (h+1)r.
 \tag{T4}
\]

This is independent of the dimension p.~Applying the generic element-label construction directly to each p-simplex would count up to p+1 vertex colors on that simplex and give r*binomial(h+p+1,p+1). The incidence construction removes that dimension-dependent degree in h. Section 5 gives a sharper exact envelope and its matching worst-case example.

\hypertarget{af-3f6acd6e5b-an-incidence-lemma-for-arbitrary-matrices}{%
\subsection{2. An incidence lemma for arbitrary matrices}\label{af-3f6acd6e5b-an-incidence-lemma-for-arbitrary-matrices}}

The useful principle is more general than simplicial homology.

Let A be an m by N matrix over F. Give row i a label set R\_i and column e a label set C\_e. Assume that every nonzero incidence satisfies

\[
R_i\subseteq C_e,
       \qquad |C_e\setminus R_i|\le d
       \quad\text{whenever } A_{ie}\ne0.
 \tag{T5}
\]

A failed label deletes a column if it lies in C\_e. There is a deterministic restriction to at most

\[
r\binom{h+d}{d}
 \tag{T6}
\]

columns preserving every surviving column span at every prefix of one fixed column order. The bound is capped by N. The size of C\_e itself need not be bounded: only its increment beyond an incident row's labels is bounded.

\hypertarget{af-3f6acd6e5b-polynomial-features}{%
\subsubsection{Polynomial features}\label{af-3f6acd6e5b-polynomial-features}}

Assign each of the L input labels a different nonzero field point x\_lambda, passing to an effective finite extension if necessary. For a label set S let

\[
g_S(T)=\prod_{\lambda\in S}(1-T/x_\lambda).
 \tag{T7}
\]

For an incidence difference Delta=C\_e minus R\_i, write

\[
g_\Delta(T)=1+\sum_{j=1}^d b_j(\Delta)T^j,
 \tag{T8}
\]

padding unused coefficients with zeros. Let Phi(Delta) be all monomials of total degree at most h in b\_1,\ldots,b\_d, including 1. Write Q=binomial(h+d,d) for the feature dimension.

For a failed-label set F\_0, the polynomial

\[
Q_{F_0}(b_1,\ldots,b_d)
   =\prod_{\lambda\in F_0}
       \left(1+\sum_{j=1}^d b_j x_\lambda^j\right)
 \tag{T9}
\]

has total degree at most h. Therefore it is a linear functional ell\_F on the feature vector. Its value is product\_\{lambda in F\_0\} g\_Delta(x\_lambda), and it vanishes precisely when Delta intersects F\_0. Failure sets are not enumerated to construct the features.

\hypertarget{af-3f6acd6e5b-retain-original-independent-rows-then-lift}{%
\subsubsection{Retain original independent rows, then lift}\label{af-3f6acd6e5b-retain-original-independent-rows-then-lift}}

Choose an index set I of r independent \textbf{original rows} of A. Do not replace these rows by arbitrary linear combinations while retaining their old label sets. Projection pi\_I is injective on the full column span of A, since all other rows are linear combinations of the selected rows.

For column e, make an r-block vector w\_e. Its block at i in I is

\[
(w_e)_i=A_{ie}\Phi(C_e\setminus R_i)
 \tag{T10}
\]

when A\_ie is nonzero, and is zero otherwise. Scan columns in the fixed order and keep exactly the columns that increase the lifted span. The retained columns form a basis in a space of dimension rQ, and every omitted column is a linear combination of earlier retained columns.

For query F\_0, multiply block i's evaluation under ell\_F by

\[
\gamma_i(F_0)=\prod_{\lambda\in F_0}g_{R_i}(x_\lambda).
 \tag{T11}
\]

Because g\_\{R\_i\}g\_\{C\_e minus R\_i\}=g\_\{C\_e\}, the resulting linear map sends

\[
w_e\longmapsto
       \left(\prod_{\lambda\in F_0}g_{C_e}(x_\lambda)\right)\pi_I(A_e).
 \tag{T12}
\]

The scalar is zero exactly on failed columns and is nonzero on every surviving column. Apply this map to each lifted prefix-span relation and divide by the protected column's nonzero scalar. It follows that every surviving original column is spanned by earlier surviving retained columns. Injectivity of pi\_I on the original column span turns this into equality in the full row space. This proves (T6) at every prefix, including prefixes inside a cost tie after one fixed ID tie break.

Ordinary matroid greedy now proves the minimum-cost-basis assertion. An original column accepted by canonical greedy after a failure cannot be an omitted column: the latter is already spanned by earlier retained survivors. Thus the entire canonical basis is contained in B, not just another basis of the same value. Signed costs and surviving rank drops cause no issue.

\hypertarget{af-3f6acd6e5b-application-to-a-filtered-simplicial-complex}{%
\subsection{3. Application to a filtered simplicial complex}\label{af-3f6acd6e5b-application-to-a-filtered-simplicial-complex}}

For a (p-1)-face tau, let R\_tau be the union of the labels of its vertices. For a p-simplex sigma, let C\_sigma be the union of its vertex labels and its own additional labels. If tau is a facet of sigma, exactly one vertex v is missing, so

\[
R_\tau\subseteq C_\sigma,
                  |C_\sigma\setminus R_\tau|\le a+b.
 \tag{T13}
\]

The signs in the simplicial boundary are just the entries A\_tau,sigma in (T10). The incidence lemma gives (T2), preserving every surviving boundary span at every filtration threshold.

The lower chain groups of H\_t-F\_0 and K\_t-F\_0 are identical in dimensions at most p-1, with the same coordinate bases and the same lower boundary maps. In degree p-1 their cycle spaces are therefore identical, and the image of partial\_p is identical by the lemma. Hence their homology quotients are identical. In lower degrees even the adjacent chain maps are unchanged. Since all these identifications use the same lower-chain coordinates, every square induced by an inclusion of states commutes. This proves equality of persistence modules.

The same naturality proves two optional extensions without increasing B:

\begin{itemize}
\tightlist
\item
  relative homology against any later specified compatible subcomplex, using H intersect that subcomplex on the retained side;
\item
  zigzag persistence for any sequence of allowed states K\_t-F\_0 whose arrows are inclusions in either direction.
\end{itemize}

These are algebraic consequences, not new claims about fast dynamic updates. For a general complex one may first pass to its p-skeleton when the goal is homology in degrees at most p-1.

\hypertarget{af-3f6acd6e5b-failures-of-higher-dimensional-faces}{%
\subsubsection{Failures of higher-dimensional faces}\label{af-3f6acd6e5b-failures-of-higher-dimensional-faces}}

Assign one label to each s-face, or use one common label for a class of s-faces, and let its failure delete those faces and all cofaces. Give a face eta the set of labels of its contained s-faces. The label sets are monotone under containment. A p-simplex has at most binomial(p,s) s-faces not contained in a specified facet, because these are the s-faces containing its missing vertex. Thus the same proof gives

\[
|B|\le r\binom{h+\binom{p}{s}}{\binom{p}{s}}.
 \tag{T14}
\]

More generally (T6) applies whenever the label assignment on faces is monotone and has incidence increment d.~This preserves answers to supplied failures; it does not optimize over the unknown worst failure set or solve the minimum homological-cut problem.

\hypertarget{af-3f6acd6e5b-uniform-construction-and-representation-details}{%
\subsection{4. Uniform construction and representation details}\label{af-3f6acd6e5b-uniform-construction-and-representation-details}}

Clip h to the number of actual labels. Compute r by exact linear algebra. If r=0 return no top columns. Otherwise compute the feature dimension Q=binomial(h+d,d) only up to the threshold ceil(N/r). If rQ\textgreater=N return every original top column. Otherwise rQ\textless N, and only then enumerate the features and the lift. All incidence differences are explicit; their polynomial coefficients and all Q monomials can be generated in polynomial time. Standard elimination uses O(N(rQ)\^{}2) field operations after feature construction and row selection. The latter costs are paid separately and are polynomial in the explicit input and generated features. This is a uniform polynomial algorithm even when h,d,p are input parameters.

Over Q, choose x\_lambda in \{1,\ldots,L\}. Coefficients in (T7)--(T9) have polynomial bit length, and exact fraction-free elimination gives polynomial bit complexity. Over a finite field of size Q\_F, if Q\_F\textgreater=L+1 no extension is needed. Otherwise choose the smallest s with Q\_F\^{}s\textgreater=L+1. Then Q\_F\textsuperscript{s\textless Q\_F(L+1)\textless(L+1)}2. Enumerating monic degree-s polynomials and testing irreducibility constructs a suitable extension deterministically in time polynomial in the explicitly present L. Original ranks are unchanged by extending scalars. Equality of spans of the original F-columns over the extension implies equality over F, by rank invariance; the homology conclusion is therefore over the original chosen coefficient field.

There is no generic-position search or unproved polynomial-identity-testing assumption. The explicit specialization may have more dependencies than a generic representation, and that does not harm the query-projection proof.

\hypertarget{af-3f6acd6e5b-a-sharp-exterior-envelope-for-one-color-per-vertex}{%
\subsection{5. A sharp exterior envelope for one color per vertex}\label{af-3f6acd6e5b-a-sharp-exterior-envelope-for-one-color-per-vertex}}

Assume now that every vertex has exactly one failable color and that b=0. Write n=\textbar V\textbar, let ell be the number of vertex colors, and put

\[
k=\min(h+1,\ell),\qquad
 E(n,p,k)=\binom n{p+1}-\binom{n-k}{p+1}.
 \tag{T15}
\]

There is a single retained set satisfying the sharper bound

\[
|B|\le\min\{N,kr,E(n,p,k)\}.
 \tag{T16}
\]

Here and below a binomial coefficient is zero when its upper argument is smaller than the nonnegative lower argument. If safe/uncolored vertices are allowed, represent them by one extra formal never-failing color when forming the color count. Otherwise the all-real-colors-failed case would be incorrectly called empty.

\hypertarget{af-3f6acd6e5b-explicit-contractions-and-the-query-map}{%
\subsubsection{Explicit contractions and the query map}\label{af-3f6acd6e5b-explicit-contractions-and-the-query-map}}

Give the colors distinct nonzero points a\_c.~Over the working field extension E, in V=E\^{}n, define covectors

\[
\ell_j(e_v)=a_{c(v)}^j,\qquad 0\le j<k.
 \tag{T17}
\]

For a p-simplex sigma, stack the contractions iota\_\{ell\_j\}(e\_sigma), with the usual oriented exterior basis. In the row block indexed by a facet tau=sigma minus v, its vector is the boundary sign times (1,a\_\{c(v)\},\ldots,a\_\{c(v)\}\^{}\{k-1\}).

For a failure set F\_0 that does not delete all colors, use P\_F(T)=product\_\{c in F\_0\}(T-a\_c), after removing colors absent from K. Its degree is at most k-1. At row tau, evaluate the feature vector against the coefficients of P\_F and multiply by product\_\{u in tau\}P\_F(a\_\{c(u)\}). The entire column maps to

\[
\left(\prod_{u\in\sigma}P_F(a_{c(u)})\right)A_\sigma.
 \tag{T18}
\]

This vanishes exactly on failed simplices. If every color fails, there are no surviving positive-dimensional simplices and the claim is immediate. The cost-ordered lifted basis therefore has the required prefix property.

\hypertarget{af-3f6acd6e5b-two-rank-bounds-for-the-same-lift}{%
\subsubsection{Two rank bounds for the same lift}\label{af-3f6acd6e5b-two-rank-bounds-for-the-same-lift}}

For each j, the corresponding weighted boundary block A\_j satisfies

\[
A_j=D_{p-1,j}^{-1} A D_{p,j},\qquad
 (D_{i,j})_{\eta\eta}=\prod_{v\in\eta}a_{c(v)}^j.
 \tag{T19}
\]

Every diagonal entry is nonzero, so rank(A\_j)=r. Stacking k blocks gives rank at most kr.

The covectors ell\_0,\ldots,ell\_\{k-1\} are independent by the Vandermonde property. Let W be their common kernel, of dimension n-k. On the complete universe Lambda\^{}\{p+1\}V, the simultaneous contraction map L satisfies

\[
\ker L=\Lambda^{p+1}W.
 \tag{T20}
\]

To see this, choose a basis dual to the k independent covectors. A wedge monomial containing any of the first k basis vectors has a nonzero, uniquely indexed coefficient in its contraction by the corresponding covector. Therefore a vector in all contraction kernels uses only the last n-k basis vectors. The converse is immediate. This argument works in every characteristic. The rank on the complete universe is E(n,p,k); restricting columns to K cannot increase it. This proves (T16).

\hypertarget{af-3f6acd6e5b-exact-worst-case-lower-bound}{%
\subsubsection{Exact worst-case lower bound}\label{af-3f6acd6e5b-exact-worst-case-lower-bound}}

Take the complete p-dimensional complex on vertices 1,\ldots,n, give every lower-dimensional face birth time 0, and give a top simplex sigma birth time min(sigma). The first k-1 vertices have distinct singleton colors. Partition the remaining vertices into the remaining ell-k+1 nonempty color classes. These assignments are possible because k\textless=ell\textless=n.

For each j\textless=k, fail the singleton colors of 1,\ldots,j-1. At time j, the surviving top simplices that have appeared are exactly the p-simplices on the remaining vertices whose minimum vertex is j (equivalently, the remaining simplices containing j). Each has a unique facet not containing j, so their boundary columns are independent. They span the boundary space of the full remaining complete complex: this is the usual cone basis. Thus omitting any of these top simplices changes H\_\{p-1\} at this allowed state. Taking the union over j\textless=k forces every top simplex meeting \{1,\ldots,k\}, namely exactly E(n,p,k) simplices.

The same construction forces those top simplices into the minimum-cost boundary basis for some failure query: assign each star a cost range below all later stars. Distinct positive perturbations within the ranges are permitted. Consequently the complete-universe envelope in (T16) is exact, for the original-top-simplex restriction model, not just order-optimal.

When p=1 this envelope is nk-binomial(k+1,2). Its ordinary individual-vertex failure version and the extremal graph example are established results of Sorkin--Steger--Zenklusen. The application here is to arbitrary correlated color classes and to higher-dimensional failure-preserving persistence. The exterior rank identity is part of established algebraic shifting/hyperconnectivity theory; it is not being claimed as a new identity.

\hypertarget{af-3f6acd6e5b-intrinsic-feature-rank-for-overlapping-vertex-labels}{%
\subsection{6. Intrinsic feature rank for overlapping vertex labels}\label{af-3f6acd6e5b-intrinsic-feature-rank-for-overlapping-vertex-labels}}

Suppose b=0, but allow up to a labels at each vertex. Write g\_\{Lambda(v)\}(T) from (T7) as 1+sum\_\{j=1\}\^{}a b\_j(v)T\^{}j, padding missing coefficients with zeros. Encode vertex v by psi(v), the vector of all degree-at-most-h monomials in these a nonconstant coefficients, including the constant monomial 1. Let

\[
s=\mathrm{rank}[\psi(v)]_{v\in V}
       \le \min\left\{n,\binom{h+a}{a}\right\}.
 \tag{T21}
\]

Each failure query is a linear functional on psi whose evaluation is zero exactly at failed vertices. Choose s independent coordinate functions spanning all coordinates of the feature vectors psi(v); every query functional then has an equivalent linear form in these s coordinates. Use the stacked contraction construction above. It gives

\[
|B|\le\min\left\{N,sr,
         \binom n{p+1}-\binom{n-s}{p+1}\right\}.
 \tag{T22}
\]

The common-kernel proof is unchanged. For the sr bound, a coordinate function w(v) may vanish on some vertices. Replace it by w(v)+z over F(z). The corresponding weighted boundary block then has rank r by the diagonal factorization. Every (r+1)-minor is the zero polynomial in z, hence is still zero at z=0, so the original block has rank at most r. Stacking the s blocks proves (T22). This is an actual feature-rank refinement; exact fixed-n sharpness is only asserted above for the one-color model.

\hypertarget{af-3f6acd6e5b-the-general-coefficient-in-t2-is-asymptotically-necessary}{%
\subsection{7. The general coefficient in (T2) is asymptotically necessary}\label{af-3f6acd6e5b-the-general-coefficient-in-t2-is-asymptotically-necessary}}

Let d=a+b and q=binomial(h+d,d). Take a universe U of h+d labels and one cone apex v\_S for each d-subset S of U. Split S into A\_S of size a and B\_S of size b. Give v\_S the vertex labels A\_S, and give every top simplex in its cone the extra labels B\_S. The base vertices have no labels.

Let t\textgreater=p and let the base contain t vertices with its complete (p-1)-skeleton. For each apex add every top simplex \{v\_S\} union tau, where tau is a p-element subset of the base; include all lower faces, and no top simplex containing two apices. The number of top simplices and the original boundary rank are

\[
M=q\binom tp,\qquad
 r=\binom tp+(q-1)\binom{t-1}{p-1}.
 \tag{T23}
\]

For completeness, write a top chain as sum\_S v\_S*C\_S, where C\_S is a (p-1)-chain on the base skeleton. Its boundary vanishes exactly when every C\_S is a cycle and sum\_S C\_S=0. The base cycle-space dimension is binomial(t-1,p), using the augmented boundary for p=1. Thus the top kernel has dimension (q-1)binomial(t-1,p), giving (T23).

Fail F\_0=U minus S. This uses exactly h labels. The top cells of cone S survive, and every other cone's top cells fail through its apex labels or its own labels. Each surviving top column has its unique base-face row. Therefore omitting any one of them decreases boundary rank and changes H\_\{p-1\}, even though other cones' lower faces may remain. Every universal restriction must retain all M top cells. Hence

\[
\frac Mr=
          \frac{q}{1+(q-1)p/t}\longrightarrow q.
 \tag{T24}
\]

This works over every coefficient field and for all a,b, including a=0 or b=0. It is a retained-original-cell lower bound with the lower skeleton fixed. For p=1 it is a simple complete bipartite graph construction, yielding the same leading coefficient for mixed overlapping vertex/edge labels.

\hypertarget{af-3f6acd6e5b-related-work-and-scope}{%
\subsection{8. Related work and scope}\label{af-3f6acd6e5b-related-work-and-scope}}

The meaningful comparison is as follows.

\begin{itemize}
\tightlist
\item
  A single no-failure boundary basis is classical. Minimum spanning higher-dimensional trees/forests and persistent skeletons are developed by Kalisnik Verovsek--Kurlin--Lesnik. Their richer persistent skeleton also addresses top-dimensional features by retaining critical faces; (T2) only promises the lower-dimensional modules with the lower skeleton fixed.
\item
  Sorkin--Steger--Zenklusen give the sharp graph bound for individual-vertex failures. Arbitrary correlated vertex-color failures are a different quantifier: one color can delete many vertices.
\item
  Kalai's exterior shifting and hyperconnectivity framework supplies the matrix/rank antecedents. The construction here uses explicit repeated color points and a query-annihilating map; it does not derandomize the generic independence problem for all algebraic-shifting matroids.
\item
  Roos Hoefgeest--Slot study adversarial robustness of persistent features and minimum homological cuts, including hardness. A certificate preserving all supplied failure queries does not solve their worst-failure search.
\item
  No identical all-color-failure persistence certificate was located in the bounded primary-source search. This is not a proof of novelty. The sharp explicit query theorem, its relation to the older exterior framework, and any earlier equivalent formulation are the principal external-review targets.
\end{itemize}

\hypertarget{af-3f6acd6e5b-primary-references}{%
\subsubsection{Primary references}\label{af-3f6acd6e5b-primary-references}}

\begin{enumerate}
\def\labelenumi{\arabic{enumi}.}
\tightlist
\item
  Sara Kalisnik Verovsek, Vitaliy Kurlin, Davorin Lesnik, \emph{A Higher-Dimensional Homologically Persistent Skeleton}. https://arxiv.org/pdf/1701.08395
\item
  Gregory B. Sorkin, Angelika Steger, Rico Zenklusen, \emph{A tight bound on the collection of edges in MSTs of induced subgraphs}, JCTB 2009; preprint 2007. https://arxiv.org/pdf/0705.2439
\item
  Gil Kalai, \emph{Hyperconnectivity of graphs}, Graphs and Combinatorics 1 (1985), 65--79. https://doi.org/10.1007/BF02582930
\item
  Lazar Guterman and Eran Nevo, \emph{Matroids arising from algebraic shifting}. https://arxiv.org/html/2512.03294v1
\item
  Pepijn Roos Hoefgeest and Lucas Slot, \emph{Robustness of Persistent Topological Features and Minimum Homological Cuts}. https://arxiv.org/html/2602.23154v1
\end{enumerate}

The theorem depends on the chosen coefficient field, retains the full lower skeleton, and uses one specified order. These hypotheses are part of its statement.
